\subsection{The tacnode and transverse-line critical curve}

\begin{statement}{P4R.16}{Proved}
In the remaining positive ruled case with uniform transverse
type $A_3$, the three-branch point of the critical curve cannot
be the union of a plane tacnode and an independent transverse
line. Its local conductor defect is at least eleven, whereas
the total defect in the contracted-section case is nine.
\end{statement}
\begin{proof}
The analytic ring calculation in P4.40 supplies actual ambient
coordinates for this curve. Its reduced ideal is
\[
 P=(xz,yz,h),\qquad h=y^2-x^4.
\]
Its three parametrizations are $(t,t^2,0)$, $(t,-t^2,0)$,
and $(0,0,t)$. The two tangent branches have intersection
multiplicity two, and their union meets the independent line
with intersection length one. The normalization sequence
therefore gives $\delta=3$. Every branch is immersive, so
its ramification divisor has coefficient zero.

We first determine the symbolic square and the actual
equation. The plane-tacnode ideal $(z,h)$ is a reduced
complete intersection. Its square is unmixed and agrees
with the square of each prime at its generic point; hence
its symbolic square is its ordinary square. Intersecting
with the square of the line ideal gives
\[
 P^{(2)}=(z,h)^2\cap(x,y)^2=P^2+(zh).
\]
Indeed an element of $(z,h)^2$ is $Az^2+Bzh+Ch^2$.
Its last two terms belong to $(x,y)^2$. The first belongs
to that ideal precisely when $A$ does, because $z$ is a
nonzerodivisor modulo $(x,y)^2$. The resulting three
generators $x^2z^2,xyz^2,y^2z^2$ are already in $P^2$.
Furthermore
\[
 xzh=(xz)h,\qquad yzh=(yz)h,\qquad
 z^2h=(yz)^2-x^2(xz)^2.
\]
The class of $zh$ modulo $P^2$ is therefore annihilated
by the ambient maximal ideal. It is nonzero, since it has
a nonzero cubic term and $P^2$ has order at least four.
Thus $P^{(2)}/P^2$ is exactly one copy of the residue field.

Let $f$ be the reduced defining function of the support.
Its generic transverse multiplicity is two, so
$f\in P^{(2)}$ and
\[
                       f=\kappa zh+f_0,\quad f_0\in P^2,
                       \quad\kappa\in\mathbb C.
\]
At $(t,t^2,0)$ with $t\ne0$, use $z,h$ as transverse
coordinates. The quadratic transverse part of $f$ has the
form
\[
 t^2a(t)z^2+(\kappa+t b(t))zh+c(t)h^2,
\]
where $a,b,c$ are holomorphic at $t=0$. The transverse
Hessian has rank one, so its determinant vanishes for
all sufficiently small nonzero $t$. Its constant term at
zero is $-\kappa^2$ (up to the harmless common convention
for the quadratic-form matrix). Hence $\kappa=0$, and
we have proved the exact inclusion $f\in P^2$.

Let $K_i$ be the saturated rank-two kernel of the derivative
of the $i$th parametrization, and let $E_i\subset K_i$ be
the span of the differentiated generators of $P$. On the
two tangent branches choose the bases
\[
                dz,\quad dy\mp2t\,dx.
\]
The generator matrices have columns
\[
 \begin{pmatrix}t&\pm t^2&0\\0&0&\pm2t^2\end{pmatrix}.
\]
Their elementary-divisor orders are $(1,2)$. On the
independent line use $dx,dy$; the matrix is
\[
                   \begin{pmatrix}t&0&0\\0&t&0\end{pmatrix},
\]
with orders $(1,1)$.

For completeness, the rank-one quadratic estimate of
P4.42 applies to these exact lattices as follows. Write
$E=t^a\mathbb C\{t\}e_1\oplus
t^b\mathbb C\{t\}e_2$, and write the saturated Hessian
line as $\ell=u e_1+v e_2$, with $(u,v)$ primitive.
Since $f\in P^2$, the tensor $t^d\ell^2$ belongs to
$\operatorname{Sym}^2E$. Therefore
\[
 \frac d2\ge\max(a-\operatorname{ord}u,
                         b-\operatorname{ord}v).
\]
Projection of $E$ to $K/\ell$ has colength
$b_{\rm proj}=\min(a+\operatorname{ord}v,
b+\operatorname{ord}u)$. Adding these two expressions
gives $d/2+b_{\rm proj}\ge a+b$; the identity follows
by subtracting the same two candidate quantities, with
the zero-component cases read directly. The local defect
formula of P4.42 consequently gives
\[
 \delta+\sum_i(d_i/2+b_i-r_i)+q
       \ge3+(1+2)+(1+2)+(1+1)=11.
\]
Here $q\ge0$ is the further descent colength, retained
rather than discarded from the definition. P4.44 contracts
the horizontal section, and P4.33 and (4.42a) give total
defect nine. The displayed lower bound contradicts that
total. This excludes the stated analytic curve germ.
\end{proof}

The symbolic-square identities, all three normalized
conormal matrices, and the rank-one determinant obstruction
are checked exactly in
\texttt{code/part4/a3\_tacnode\_line.py}.
